\chapter{CSF IgG Index}

\section*{What Is This Test?}
The CSF IgG index (also called the IgG index or Link index) quantifies intrathecal synthesis of immunoglobulin G (IgG). Normally, IgG in CSF derives primarily from plasma, with only a small fraction produced centrally by B cells in the central nervous system. In chronic CNS inflammatory diseases, especially multiple sclerosis (MS), intrathecal plasma cells synthesize IgG locally, increasing CSF IgG out of proportion to albumin. The index compares CSF and serum IgG/albumin ratios and is the most widely used qualitative marker of intrathecal humoral immunity.

\section*{Alternative Names}
CSF IgG Index; IgG Index; Link Index; CSF Immunoglobulin G Index; CSF/Serum IgG-Albumin Ratio

\section*{Principle}
Albumin is synthesized exclusively in the liver; it does not cross the intact blood-CSF barrier under normal conditions, so CSF albumin increases passively with barrier disruption. IgG is also plasma-derived but also locally synthesized in inflammatory CNS disease. The index is calculated as:
\[
IgG Index = \frac{[IgG\textsubscript{\text{CSF}}]/[IgG\textsubscript{\text{serum}}]}{[Albumin\textsubscript{\text{CSF}}]/[Albumin\textsubscript{\text{serum}}]}
\]
A value $>$ 0.7 indicates intrathecal IgG synthesis. Paired serum and CSF samples are required on the same day. Quantitation by immunonephelometry or immunoturbidimetry on automated analyzers \cite{bonning2020cerebrospinal}.

\section*{History}
The IgG index was introduced by Tibbling and Link in 1977 as a refinement over qualitative CSF IgG measurements \cite{link1977principles}. It remains the principal quantitative marker of intrathecal IgG synthesis, complemented by qualitative oligoclonal bands (OCBs) detection by isoelectric focusing \cite{freedman2005recommended}.

\section*{Why This Test Is Ordered}
\begin{itemize}
  \item Diagnostic workup of suspected multiple sclerosis (2017 McDonald criteria) \cite{thompson2018diagnosis}
  \item Evaluation of chronic CNS inflammatory or infectious diseases: neuroborreliosis, neurosyphilis, subacute sclerosing panencephalitis (SSPE), viral encephalitis
  \item Investigation of unexplained CNS demyelination, transverse myelitis, optic neuritis
  \item Monitoring Intrathecal immune activity in immune reconstitution inflammatory syndromes (IRIS)
\end{itemize}

\section*{Is This a Primary or Secondary Test}
Secondary test; complements MRI and oligoclonal band testing in MS diagnostic workup.

\section*{How This Test Is Performed}
Cerebrospinal fluid (2--5 mL) is collected via lumbar puncture in sterile tubes. A concurrent serum sample (5 mL in serum separator tube) must be drawn on the same day. Both are assayed for IgG and albumin; the index is calculated. Oligoclonal bands should be ordered simultaneously for full evaluation.

\section*{What Preparation Is Needed}
No specific prep besides patient consent for lumbar puncture. Imaging (brain MRI) may precede LP; in MS workup, perform MRI before LP to minimize post-LP headache appearing to be neurologic symptom.

\section*{Contraindications and Precautions}
Lumbar puncture contraindications: anticoagulation or coagulopathy, overlying soft tissue infection, suspected raised intracranial pressure with mass effect. Traumatic tap may elevate CSF total protein and IgG; in such cases, correction formulas (e.g., using RBC counts) can be used.

\section*{Risks}
Post-LP headache (10--20\%); local back pain; rare: infection, bleeding, brain herniation if raised ICP.

\section*{What Happens After the Test}
Results reported within 1--3 days. An elevated IgG index in an appropriate clinical context supports MS diagnosis; interpretation is enhanced by oligoclonal bands, MRI lesions, symptoms, and exclusion of mimics.

\section*{Understanding Results}

\subsection*{Below Normal}
A low or normal index rules out significant intrathecal IgG synthesis but does not exclude all cases of early MS; oligoclonal bands can be present with a normal IgG index.

\subsection*{Above Normal}
\begin{itemize}
  \item \textbf{Multiple sclerosis:} IgG index $>$ 0.7 in approximately 70--90\%; combined with oligoclonal bands are 2017 McDonald criteria
  \item \textbf{Neuroborreliosis (Lyme neuroborreliosis):} Often elevated; antibody index $>$ 1.3 confirms
  \item \textbf{Neurosyphilis:} CSFVDRL and FTA-ABS positive; IgG index often elevated
  \item \textbf{Subacute sclerosing panencephalitis (SSPE):} Strikingly elevated CSF IgG and measles antibody
  \item \textbf{Viral encephalitis}: variable
  \item \textbf{Guillain-Barré syndrome, CIDP}: albuminocytologic dissociation; index usually normal because process peripheral, not central
\end{itemize}

\section*{Normal Results}
\begin{itemize}
  \item \textbf{CSF IgG index:} 0.3--0.7
  \item \textbf{Intrathecal synthesis:} $>$ 0.7
  \item \textbf{CSF IgG:} 5--60 mg/dL (approximate, depends on age and barrier function)
  \item \textbf{CSF Albumin:} 13--35 mg/dL
  \item \textbf{CSF/Serum Albumin Ratio (Q Alb):} 2--8 $\times$ 10$^{-3}$ (barrier integrity marker)
\end{itemize}

\section*{Follow-up Tests}
\begin{itemize}
  \item \textbf{Oligoclonal bands (CSF and serum):} Qualitative confirmation of intrathecal IgG; $>$ 2 unique CSF bands not present in serum is positive
  \item \textbf{Brain and spinal cord MRI (with gadolinium):} MS dissemination in space and time
  \item \textbf{Visual evoked potentials:} Optic neuritis assessment
  \item \textbf{Lyme (antibody index) and syphilis serologies (CSF VDRL, serum RPR):} For exclusion
  \item \textbf{CSF cell count, protein, glucose, culture:} Basic CSF evaluation
\end{itemize}

\section*{References}
\bibliographystyle{unsrtnat}
\bibliography{references}

\clearpage